% Unit 119 terjemahan Bahasa Indonesia dari chapter8.tex baris 2000--2107.
% Sumber: Methods of Algebra, Volume 2, commit
% 9a5803ff77dd3257484cb177f851a73770a59dd3 (CC BY 4.0).
% stable-unit-id: o014.aljabr2.chapter8.exercises-a
% Cakupan: paruh pertama latihan Bab 8, dari pembukaan Exercises hingga
% definisi objek L-proyektif dan T-injektif; Unit 120 melanjutkan baris 2108.

\hypertarget{o014.aljabr2.chapter8.exercises-a}{ }
% segment-id: o014.aljabr2.chapter8.exercises-a.h001
\begin{Exercises}
% segment-id: o014.aljabr2.chapter8.exercises-a.ex001
	\item Buktikan bahwa, untuk setiap kategori $\mathcal{C}$, konstruksi
	$C\mapsto\mathrm{const}(C)$ dari Contoh~\ref{eg:const-simplicial}
	memberikan funktor penuh dan setia
	$\mathcal{C}\to\cate{s}\mathcal{C}$.

% segment-id: o014.aljabr2.chapter8.exercises-a.ex002
	\item Misalkan $X$ himpunan simplisial dan $x\in X_n$. Buktikan bahwa
	terdapat barisan indeks $j_1<\cdots<j_h$ dan suatu simpleks nondegenerat
	$y$ sedemikian sehingga $x=s_{j_h}\cdots s_{j_1}(y)$, dan bahwa $y$ unik.

% segment-id: o014.aljabr2.chapter8.exercises-a.hi001
	\begin{hint}
		Kesulitannya terletak pada keunikan. Misalkan
		$Sy=x=S'y'$, dengan $S$ dan $S'$ keduanya berbentuk seperti dalam
		pernyataan. Kita dapat memilih komposisi $D$ dari serangkaian morfisme muka
		sedemikian sehingga $y=DSy=DS'y'$. Ubah $DS'$ ke bentuk
		$\widetilde{S}'\widetilde{D}$ (makna notasinya jelas). Karena $y$
		nondegenerat, harus berlaku $\widetilde{S}'=\identity$, sehingga
		$y'=\widetilde{D}y$. Gunakan simetri untuk menyimpulkan
		$y=y'$ dan $\widetilde{D}=\identity$.
	\end{hint}

% segment-id: o014.aljabr2.chapter8.exercises-a.ex003
	\item Tunjukkan bahwa terdapat diagram berikut dalam $\cate{sSet}$, dengan
	bagian atas maupun bawah yang komutatif, untuk
	$0\leq i<j\leq n$:
% segment-id: o014.aljabr2.chapter8.exercises-a.q001
	\[\begin{tikzcd}
		\Delta^{n-2} \arrow[d, "{\iota_{i < j}}"'] \arrow[r, "{\mathrm{d}^{j-1}}"] & \Delta^{n-1} \arrow[d, "\iota_i"] \arrow[rd, "{\mathrm{d}^i}"] & \\
		\bigsqcup_{0 \leq i' < j' \leq n} \Delta^{n-2} \arrow[shift left, r] \arrow[shift right, r] & \bigsqcup_{i' = 0}^n \Delta^{n-1} \arrow[r] & \partial \Delta^n \\
		\Delta^{n-2} \arrow[u, "{\iota_{i < j}}"] \arrow[r, "{\mathrm{d}^i}"'] & \Delta^{n-1} \arrow[u, "\iota_j"'] \arrow[ru, "{\mathrm{d}^j}"'] &
	\end{tikzcd}\]
	Di sini, $\bigsqcup$ diambil suku demi suku dalam
	$\cate{sSet}=\cate{Set}^{\simpDelta^{\opp}}$, sedangkan
	$\iota_{i<j}$ (atau $\iota_i$) menyatakan embedding ke suku ke-$i<j$
	(atau suku ke-$i$). Semua anak panah pada baris tengah dicirikan oleh
	komutativitas tersebut. Tunjukkan lebih lanjut bahwa baris tengah itu
	merupakan koekualiser.

	Secara serupa, tunjukkan bahwa untuk setiap $0\leq k\leq n$ terdapat pula
% segment-id: o014.aljabr2.chapter8.exercises-a.q002
	\[\begin{tikzcd}
		\Delta^{n-2} \arrow[d, "{\iota_{i < j}}"'] \arrow[r, "{\mathrm{d}^{j-1}}"] & \Delta^{n-1} \arrow[d, "\iota_i"] \arrow[rd, "{\mathrm{d}^i}"] & \\
		\bigsqcup_{\substack{0 \leq i' < j' \leq n \\ i', j' \neq k}} \Delta^{n-2} \arrow[shift left, r] \arrow[shift right, r] & \bigsqcup_{\substack{0 \leq i' \leq n \\ i' \neq k}} \Delta^{n-1} \arrow[r] & \Lambda^n_k \\
		\Delta^{n-2} \arrow[u, "{\iota_{i < j}}"] \arrow[r, "{\mathrm{d}^i}"'] & \Delta^{n-1} \arrow[u, "\iota_j"'] \arrow[ru, "{\mathrm{d}^j}"'] &
	\end{tikzcd}\]
	dengan baris tengah merupakan koekualiser dan $i,j\neq k$. Cobalah
	jelaskan makna topologis koekualiser tersebut.

% segment-id: o014.aljabr2.chapter8.exercises-a.ex004
	\item Verifikasi sifat-sifat berikut dari operasi join $\star$ dalam
	Definisi~\ref{def:simplicial-join}.
	\begin{enumerate}[(i)]
% segment-id: o014.aljabr2.chapter8.exercises-a.i001
		\item Himpunan simplisial secara alami membentuk kategori monoidal
		terhadap $\star$, dengan himpunan simplisial kosong yang bernilai konstan
		$\emptyset$ sebagai objek satuannya.
% segment-id: o014.aljabr2.chapter8.exercises-a.i002
		\item Dualitas pembalikan urutan dari
		Catatan~\ref{rem:simpDelta-order-reversal} memberikan automorfisme
		$\cate{sSet}$, yang untuk sementara ditulis $X\mapsto\hat{X}$. Tunjukkan
		bahwa $(X\star Y)^{\wedge}\simeq\hat{Y}\star\hat{X}$.
% segment-id: o014.aljabr2.chapter8.exercises-a.i003
		\item Verifikasi
		$\Delta^p\star\Delta^q\simeq\Delta^{p+q+1}$.
	\end{enumerate}

% segment-id: o014.aljabr2.chapter8.exercises-a.ex005
	\item Tunjukkan bahwa pengambilan dual pembalikan urutan dari suatu
	himpunan simplisial, hingga homeomorfisme kanonik, tidak mengubah realisasi
	geometrisnya.
% segment-id: o014.aljabr2.chapter8.exercises-a.hi002
	\begin{hint}
		Mulailah dengan kasus $\Delta^n$.
	\end{hint}

% segment-id: o014.aljabr2.chapter8.exercises-a.ex006
	\item Verifikasi pernyataan Contoh~\ref{eg:cone-realization}.

% segment-id: o014.aljabr2.chapter8.exercises-a.ex007
	\item Misalkan $\mathcal{C}$ dan $\mathcal{C}'$ kategori. Definisikan
	kategori $\mathcal{C}\star\mathcal{C}'$ sebagai berikut:
% segment-id: o014.aljabr2.chapter8.exercises-a.q003
	\begin{align*}
		\Obj(\mathcal{C} \star \mathcal{C}') & := \Obj(\mathcal{C}) \sqcup \Obj(\mathcal{C}'), \\
		\Hom_{\mathcal{C} \star \mathcal{C}'}(X, Y) & := \begin{cases}
			\Hom_{\mathcal{C}}(X, Y), & X, Y \in \Obj(\mathcal{C}) \\
			\Hom_{\mathcal{C}'}(X, Y), & X, Y \in \Obj(\mathcal{C}'), \\
			\text{himpunan satu titik}, & X \in \Obj(\mathcal{C}), \; Y \in \Obj(\mathcal{C}'), \\
			\emptyset, & X \in \Obj(\mathcal{C}'), \; Y \in \Obj(\mathcal{C}).
		\end{cases}
	\end{align*}
	Komposisi morfisme dan morfisme identitas didefinisikan dengan cara yang
	langsung. Buktikan bahwa operasi join $\star$ pada himpunan simplisial dan
	funktor saraf dari Contoh~\ref{eg:nerve-cat} berhubungan melalui
% segment-id: o014.aljabr2.chapter8.exercises-a.q004
	\[ \mathrm{N}(\mathcal{C}) \star \mathrm{N}(\mathcal{C}') \simeq \mathrm{N}(\mathcal{C} \star \mathcal{C}'). \]

% segment-id: o014.aljabr2.chapter8.exercises-a.ex008
	\item Berikan bukti yang lebih rinci untuk
	Persamaan~\eqref{eqn:shuffle-sign-0} dan
	\eqref{eqn:shuffle-sign-1}.

% segment-id: o014.aljabr2.chapter8.exercises-a.ex009
	\item Untuk objek simplisial $X$ dalam sebarang kategori, definisikan
	objek simplisial yang disebut pergeseran, $\mathrm{Déc}_0X$ dan
	$\mathrm{Déc}^0X$, melalui
	$(\mathrm{Déc}_0X)_n=(\mathrm{Déc}^0X)_n=X_{n+1}$ dan
% segment-id: o014.aljabr2.chapter8.exercises-a.q005
	\begin{gather*}
		d^{\mathrm{Déc}_0 X, n}_i = d^{X, n+1}_i, \quad s^{\mathrm{Déc}_0 X, n}_j = s^{X, n+1}_j, \\
		d^{\mathrm{Déc}^0 X, n}_i = d^{X, n+1}_{i+1}, \quad
		s^{\mathrm{Déc}^0 X, n}_j = s^{X, n+1}_{j+1}.
	\end{gather*}
	Tunjukkan bahwa keduanya terdefinisi dengan baik dan saling berhubungan
	melalui dualitas pembalikan urutan
	(Catatan~\ref{rem:simpDelta-order-reversal}). Tunjukkan pula bahwa untuk
	kategori aditif berlaku
	$\mathrm{C}(\mathrm{Déc}^0X)=\mathrm{C}(X)[1]$; lihat
	Catatan~\ref{rem:chain-cplx-translation}.
	\index{objek simplisial!pergeseran@pergeseran (décalage)}
	\index[sym1]{Dec@$\mathrm{Déc}_0 X, \mathrm{Déc}^0 X$}

% segment-id: o014.aljabr2.chapter8.exercises-a.ex010
	\item Dengan meneruskan latihan sebelumnya, tunjukkan bahwa
	$d_0:X_1\to X_0$ (atau $d_1:X_1\to X_0$) memberikan augmentasi pada
	$\mathrm{Déc}_0X$ (atau $\mathrm{Déc}^0X$), dengan $X_0$ sebagai suku
	derajat $-1$. Tunjukkan lebih lanjut bahwa setelah augmentasi,
	$\mathrm{Déc}_0X$ dapat dikontraksi dari kanan dan $\mathrm{Déc}^0X$
	dapat dikontraksi dari kiri; lihat Definisi~\ref{def:contractible-aug}.
% segment-id: o014.aljabr2.chapter8.exercises-a.hi003
	\begin{hint}
		Cukup tangani $\mathrm{Déc}_0X$. Ambil
		$k_{n-1}:=s_n:X_n\to X_{n+1}$.
	\end{hint}

% segment-id: o014.aljabr2.chapter8.exercises-a.ex011
	\item Buktikan bahwa kategori kecil $\mathcal{C}$ merupakan grupoid
	(yakni semua morfismenya invertibel) jika dan hanya jika
	$\mathrm{N}\mathcal{C}$ memenuhi syarat berikut: untuk setiap
	$n\in\Z_{\geq1}$ dan $0\leq i\leq n$, setiap morfisme
	$\Lambda^n_i\to\mathrm{N}\mathcal{C}$ dapat diperluas menjadi
	$\Delta^n\to\mathrm{N}\mathcal{C}$, tanpa menuntut keunikan. Himpunan
	simplisial yang memenuhi syarat perluasan ini disebut kompleks Kan.

% segment-id: o014.aljabr2.chapter8.exercises-a.ex012
	\item Lanjutkan pembahasan Contoh~\ref{eg:classifying-space-std-cplx}.
	Tunjukkan bahwa kompleks ternormalisasi $\overline{\mathsf{L}}$, setelah
	dibalik menjadi kompleks rantai, isomorfik dengan
	$\mathrm{N}(\Bbbk\mathrm{E}\Gamma)$.
% segment-id: o014.aljabr2.chapter8.exercises-a.hi004
	\begin{hint}
		Pertama, gunakan Proposisi~\ref{prop:Dold-Kan-v} untuk mengidentifikasi
		$\mathrm{N}(X)$ dengan suatu hasil bagi dari $\mathrm{C}(X)$.
	\end{hint}

% segment-id: o014.aljabr2.chapter8.exercises-a.ex013
	\item Misalkan $\Gamma$ monoid. Tuliskan $\cated{Set}\Gamma$ untuk kategori
	semua himpunan kecil dengan aksi kanan $\Gamma$.
	\begin{enumerate}[(i)]
% segment-id: o014.aljabr2.chapter8.exercises-a.i004
		\item Deskripsikan secara eksplisit pasangan adjoin bebas--pelupa
% segment-id: o014.aljabr2.chapter8.exercises-a.q006
		\[\begin{tikzcd}
			\mathrm{F}: \cate{Set} \arrow[shift left, r] & \cated{Set}\Gamma \arrow[shift left, l]: U
		\end{tikzcd}\]
		beserta morfisme unit dan kounitnya. Tunjukkan bahwa monad $T$ yang
		bersesuaian memberikan ekuivalensi
		$\cated{Set}\Gamma\simeq\cate{Set}^T$.

% segment-id: o014.aljabr2.chapter8.exercises-a.hi005
		\begin{hint}
			Funktor $\mathrm{F}$ memetakan himpunan $X$ ke $X\times\Gamma$,
			dengan $\Gamma$ bekerja melalui perkalian kanan pada komponen kedua.
			Morfisme unit adalah $x\mapsto(x,1)$, sedangkan morfisme kounit
			adalah peta aksi.
		\end{hint}
% segment-id: o014.aljabr2.chapter8.exercises-a.i005
		\item Untuk setiap himpunan-$\Gamma$ kanan $X$, deskripsikan objek
		simplisial $\mathrm{Bar}(X)$ yang bersesuaian dalam
		$\cated{Set}\Gamma$; lihat Contoh~\ref{eg:bar-adjunction} dan
		\ref{eg:bar-comonad}.
% segment-id: o014.aljabr2.chapter8.exercises-a.i006
		\item Ambil himpunan satu titik $X=\{\mathrm{pt}\}$ dengan aksi kanan
		$\Gamma$ trivial. Tunjukkan bahwa objek simplisial yang bersesuaian
		isomorfik dengan $\mathrm{E}\Gamma$ dari
		Contoh~\ref{eg:classifying-space}; objek tersebut menjadi objek
		simplisial dalam $\cated{Set}\Gamma$ melalui aksi kanan $\Gamma$ pada
		setiap $(\mathrm{E}\Gamma)_n$.
	\end{enumerate}

% segment-id: o014.aljabr2.chapter8.exercises-a.ex014
	\item Tetapkan grup $\Gamma$. Dalam Contoh~\ref{eg:HH-comonad}, ambil
	aljabar grup $R=\Bbbk[\Gamma]$ dan ambil $M$ sebagai modul-$\Gamma$
	trivial $\Bbbk$; lihat Contoh~\ref{eg:trivial-G-mod}. Tunjukkan bahwa
	$\epsilon_{\Bbbk}:\mathrm{C}(\mathrm{Bar}(\Bbbk))\to\Bbbk$ yang
	bersesuaian isomorfik dengan resolusi
	$\mathsf{L}=(\mathsf{L}_n,\partial'_n)_n\to\Bbbk$ dari
	Proposisi~\ref{prop:trivial-mod-resolution}.
% segment-id: o014.aljabr2.chapter8.exercises-a.hi006
	\begin{hint}
		Gunakan \eqref{eqn:L-basis-partial}.
	\end{hint}

% segment-id: o014.aljabr2.chapter8.exercises-a.ex015
	\item Ambil gelanggang komutatif $\Bbbk$ dan
	$\mathcal{A}=\Bbbk\dcate{Mod}$, dengan
	$\otimes=\otimes_{\Bbbk}$. Misalkan $\Gamma$ grup. Tinjau objek
	simplisial $\Bbbk\mathrm{E}\Gamma$ dari
	Contoh~\ref{eg:classifying-contractible}. Terapkan
	Definisi~\ref{def:Eilenberg-Zilber} tentang morfisme Alexander--Whitney
	pada $X:=\Bbbk\mathrm{E}\Gamma\boxtimes\Bbbk\mathrm{E}\Gamma$,
	lalu tinjau morfisme
% segment-id: o014.aljabr2.chapter8.exercises-a.q007
	\[ \mathrm{C}(\Bbbk\mathrm{E}\Gamma) \to \mathrm{C}(\underbracket{\Bbbk\mathrm{E}\Gamma \otimes \Bbbk\mathrm{E}\Gamma}_{\simeq \Bbbk(\mathrm{E}\Gamma \times \mathrm{E}\Gamma)}) \xrightarrow{\overline{\mathrm{AW}}} \mathrm{C}(\Bbbk\mathrm{E}\Gamma) \otimes \mathrm{C}(\Bbbk\mathrm{E}\Gamma). \]
	Anak panah pertama diinduksi oleh embedding diagonal himpunan simplisial
	$\mathrm{E}\Gamma\to\mathrm{E}\Gamma\times\mathrm{E}\Gamma$. Mengingat
	Contoh~\ref{eg:classifying-space-std-cplx}, komposisinya juga dapat
	dipandang sebagai morfisme kompleks rantai
% segment-id: o014.aljabr2.chapter8.exercises-a.q008
	\[ \mathsf{L} \to \mathsf{L} \otimes \mathsf{L}. \]
	Gunakan ini untuk menafsirkan morfisme
	$\Delta:\mathsf{L}\to\mathsf{L}\otimes\mathsf{L}$ yang digunakan dalam
	definisi hasil kali cawan pada kohomologi grup; lihat
	Lema~\ref{prop:cup-resolution}~(ii).

% segment-id: o014.aljabr2.chapter8.exercises-a.hi007
	\begin{hint}
		Dengan menyubstitusikan definisi $\overline{\mathrm{AW}}$, terlihat
		bahwa komponen $(p,q)$ dari morfisme komposit berbeda dengan yang ada
		dalam Lema~\ref{prop:cup-resolution}~(ii) sebesar faktor $(-1)^{pq}$.
		Perbedaan ini muncul karena pembahasan di sana menggunakan
		$(\mathsf{L}_n,\partial_n)_n$, sedangkan
		Contoh~\ref{eg:classifying-space-std-cplx} menggunakan
		$(\mathsf{L}_n,\partial'_n)_n$. Pada tingkat himpunan simplisial
		$\mathrm{E}\Gamma$, keduanya berbeda dengan suatu pembalikan urutan
		(Catatan~\ref{rem:simpDelta-order-reversal}). Perhatikan bahwa dualitas
		pembalikan urutan mempertukarkan peran $\lambda^n_p$ dan $\rho^n_q$
		dalam $\overline{\mathrm{AW}}$, atau dengan kata lain mempertukarkan
		kedua salinan $\mathrm{C}(\Bbbk\mathrm{E}\Gamma)$.
	\end{hint}

% segment-id: o014.aljabr2.chapter8.exercises-a.ex016
	\item Tinjau komonad $(L,\delta,\epsilon)$ pada kategori $\mathcal{D}$.
	Untuk setiap $\mathcal{M}\in\Obj(\mathcal{D})$, jika
	$\epsilon_{\mathcal{M}}:L\mathcal{M}\to\mathcal{M}$ mempunyai invers
	kanan $f:\mathcal{M}\to L\mathcal{M}$, maka $\mathcal{M}$ disebut
	\emph{objek $L$-proyektif}. Secara dual, dari monad $(T,\mu,\eta)$ pada
	kategori $\mathcal{C}$ kita mendefinisikan
	\emph{objek $T$-injektif} dalam $\mathcal{C}$.
